\documentclass[12pt]{article}
\usepackage[a4paper,margin=1in]{geometry}
\usepackage[T1]{fontenc}
\usepackage[utf8]{inputenc}
\usepackage{lmodern}
\usepackage{setspace}
\usepackage{enumitem}
\usepackage{xcolor}

\setstretch{1.15}
\setlength{\parindent}{0pt}
\setlength{\parskip}{0.6em}

\begin{document}

\begin{center}
    {\LARGE \textbf{Presentation Script}}\\[0.5em]
    {\large Why Study Economics at UNBC? Understanding Real-World Issues}
\end{center}

\section*{Slide 1 --- Title slide}
\textbf{Why Study Economics at UNBC? Understanding Real-World Issues}

``Good afternoon everyone. [smile] Thank you so much for being here.

My name is Dr. Muhebullah Karimzada, and today I want to talk to you about a simple but very important question: \textbf{why study economics at UNBC?} [pause]

Now, when some people hear the word \textit{economics}, they immediately think of inflation, stock markets, or very complicated graphs. [small smile] And some people think of a class that may or may not include suffering. [light laughter]

But economics is actually much more interesting than that. It is about how people make choices, how the world works, and how we can better understand the challenges we face in everyday life. [pause]

So today, I want to make economics feel approachable, relevant, and hopefully even a little fun.'' [smile]

\section*{Slide 2 --- What is Economics?}

``Let me start with the most basic question: \textbf{what is economics?} [pause]

This is an important question because many people use the word, but not everyone is quite sure what it really means.

A simple answer is this: economics is the study of how people, businesses, and governments make choices when resources are limited. [pause]

That may sound formal, but the idea is actually very human and very practical.''

\newpage
\section*{Slide 3 --- What is economics?}

``In economics, one of the key ideas is \textbf{scarcity}. [pause]

Scarcity simply means that time, money, energy, and opportunities are limited. We all want many things, but we cannot have everything at once. [pause]

Because of that, every choice involves a \textbf{trade-off}. If I spend time on one thing, I give up time for something else. If a government spends more money on one program, it may have less for another. [pause]

And this is where economics becomes useful. Economics helps us think clearly about those trade-offs. [pause]

Take a student, for example. A student may want more sleep, higher grades, a part-time job, time with friends, and time to relax. [smile] That sounds reasonable. [pause] The problem is that there are only 24 hours in a day. [light laughter]

That is economics.''

\section*{Slide 4 --- Economics is everywhere}

``One of the biggest misconceptions about economics is that it is only about money. It is not. [pause]

Economics is everywhere. It is at home when families make grocery budgets, decide between saving and spending, or manage rent, mortgage payments, and monthly bills. [pause]

It is at school when students choose courses, manage their time, or decide whether to work while studying. [pause]

And it is in society when we think about inflation, wages, housing affordability, taxes, public services, and climate policy. [pause]

In other words, economics is what you are already doing --- even if you never call it economics.'' [smile]
\newpage
\section*{Slide 5 --- A few real-life examples}

``Let me make this even more concrete with a few everyday examples. [pause]

Think about a coffee shop line. Why is there such a long line sometimes? Because demand is high and supply is limited. [pause]

Think about concert tickets. When there are only so many seats and thousands of people want them, prices rise --- especially in resale markets. [pause]

Think about grocery shopping. When prices rise, families often switch to cheaper alternatives. That is economics too. [pause]

And think about choosing a university degree. That choice also has an \textbf{opportunity cost}. [pause]

Opportunity cost means that the real cost of a choice is the value of the next best alternative you give up. [pause]

So if you spend three hours scrolling on your phone instead of studying, economics politely asks: \textit{what did that really cost you?}'' [smile] [light laughter]

\section*{Slide 6 --- The big questions economics helps explain}

``Economics also helps us understand some of the biggest questions in society. [pause]

Why do prices rise? That is inflation.\\
Why can jobs be hard to find? That is unemployment.\\
Why is housing so expensive? That involves supply, demand, and policy.\\
Why do some people or places do better than others? That brings us to growth and inequality. [pause]

And how should society respond to pollution and climate change? That is where environmental economics becomes important. [pause]

One of the most valuable lessons in economics is that we do not stop at the first effect. We ask: \textbf{what happens next?} [pause]

That habit of thinking is useful not only in economics, but in life more generally.''
\newpage
\section*{Slide 7 --- Economics is about people, not just numbers}

``Sometimes people think economics is only about numbers, formulas, and graphs. [pause]

But economics is really about people. [pause]

It is about families deciding how to spend their income.\\
It is about students choosing their future.\\
It is about workers looking for secure jobs.\\
It is about governments trying to use limited budgets wisely.\\
And it is about communities trying to improve quality of life. [pause]

So good economics is deeply human. It studies well-being, fairness, incentives, and the consequences of choices. [pause]

The numbers matter, yes --- but they matter because people matter.''

\section*{Slide 8 --- Behavioural economics}

``Now let me move to one of the most interesting areas of economics: \textbf{behavioural economics}. [pause]

Traditional economics often assumes that people are rational and make careful decisions. Behavioural economics says: [pause, smile] \textit{that is a nice idea, but have you met actual people?} [laughter]

Behavioural economics studies how people really behave. It looks at procrastination, habits, overconfidence, peer effects, and self-control problems. [pause]

For example:\\
Why do we buy things we do not need?\\
Why do people delay saving for retirement?\\
Why is `starting tomorrow' such a popular life strategy? [smile]

Behavioural economics reminds us that humans are not robots. [pause] Sometimes we are not even very good spreadsheets.'' [light laughter]
\newpage
\section*{Slide 9 --- Neuroeconomics}

``Another fascinating field is \textbf{neuroeconomics}. [pause]

This area connects economics with psychology and neuroscience. It studies what happens in the brain when people make decisions. [pause]

It helps us better understand risk-taking, rewards, trust, addiction, and impulsive choices. [pause]

For instance, why do people take risks?\\
How do emotions affect financial decisions?\\
Why does immediate reward often beat long-term benefit? [pause]

And no --- neuroeconomics does not read minds. [smile] Economics still struggles to predict what a family wants for dinner when everyone says, \textit{anything is fine.}'' [laughter]

\section*{Slide 10 --- Other fascinating fields in economics}

``And these are just two examples. Economics is actually a very broad and modern field. [pause]

There is \textbf{labour economics}, which studies jobs, wages, skills, and unemployment.\\
There is \textbf{health economics}, which looks at healthcare costs and policy.\\
There is \textbf{environmental economics}, focused on climate, pollution, and sustainability. [pause]

There is also \textbf{development economics}, which studies poverty, education, and growth.\\
\textbf{Financial economics}, which deals with banking, investing, and risk.\\
And \textbf{public economics}, which looks at taxes, spending, and government policy. [pause]

So the big takeaway here is that economics is not narrow. It is broad, modern, and connected to many of the major issues in the world.''
\newpage
\section*{Slide 11 --- Why economics matters now more than ever}

``Economics matters now more than ever because we are living through major economic changes. [pause]

We see rising living costs and inflation.\\
We see housing and affordability pressures.\\
We see artificial intelligence changing labour markets.\\
We see global trade disruptions, supply chain problems, and uncertainty.\\
And we see climate and sustainability challenges. [pause]

Economics gives students a framework for understanding all of this --- instead of simply reacting to headlines. [pause]

That is one of the real strengths of studying economics: it helps you make sense of the world.''

\section*{Slide 12 --- What skills does economics build?}

``Now let us talk about skills. [pause]

Economics builds critical thinking.\\
It builds problem solving.\\
It builds data interpretation.\\
It teaches decision-making under uncertainty.\\
And it strengthens writing and communication. [pause]

It also helps students understand incentives and unintended consequences --- which is a very powerful way of thinking. [pause]

These are not only economist skills. They are career skills and life skills. [pause]

Even students who do not become economists benefit from learning how to analyze problems clearly and make evidence-based decisions.''
\newpage
\section*{Slide 13 --- Why students should consider economics}

``So why should students consider economics? [pause]

Because it is practical.\\
Because it is relevant.\\
Because it combines numbers, ideas, and real-world issues.\\
And because it opens doors to many careers. [pause]

Economics works well on its own, but it also combines very well with business, political science, data science, environmental studies, and many other fields. [pause]

Students with economics training can go into government and public policy, banking and finance, business and consulting, research and analytics, law, education, and graduate study. [pause]

So economics does not close doors. It opens them.''

\section*{Slide 14 --- A message for parents and students}

``Let me say a brief word directly to both students and parents. [pause]

For students: if you are curious about how the world works, and if you enjoy asking \textit{why}, economics is worth exploring. [pause]

For parents: economics helps develop adaptable, analytical, employable graduates with skills valued across many sectors. [pause]

And perhaps most importantly --- after one economics course, students may finally understand why \textit{we cannot buy everything} is not just parenting. [smile] It is scarcity.'' [laughter]

\section*{Slide 15 --- Final takeaway}

``So let me close with this main idea. [pause]

Economics is everywhere.\\
It helps us understand choices, incentives, markets, policy, and human behaviour. [pause]

It is relevant to everyday life.\\
It helps explain major social and economic challenges.\\
It includes exciting fields like behavioural economics and neuroeconomics.\\
And it builds valuable skills with strong career opportunities. [pause]

Economics is not just a subject. [pause] It is a way of understanding the world.''

\section*{Slide 16 --- Thank you / Questions}

``Thank you very much for your time and attention. [smile]

I hope this has helped make economics feel a little more familiar, a little more practical, and maybe even a little more exciting. [pause]

I would be very happy to take your questions.''

\section*{Delivery tips}

A few things will make this sound especially strong in person:

Pause after your opening question on each major slide. It gives the audience time to think.

Smile when you deliver the lighter lines, especially the jokes about spreadsheets, dinner choices, and scarcity. Those are gentle and audience-friendly.

Do not rush the ``economics is about people'' slide. That is one of the strongest emotional points in the presentation.


\end{document}